% Vietnamese translation for OI Public Library
% Source-PDF: IOI中国国家候选队论文/国家集训队2017论文集.pdf
\documentclass[11pt]{article}
\usepackage{oipl-vn}

\newcommand{\Oh}{\mathcal{O}}
\newcommand{\Base}{\operatorname{Base}}
\newcommand{\Res}{\operatorname{res}}
\newcommand{\Pro}{\operatorname{pro}}
\newcommand{\num}{\operatorname{num}}
\newcommand{\last}{\operatorname{last}}
\newcommand{\append}{\operatorname{append}}
\newcommand{\qedmark}{\hfill\rule{0.7ex}{0.7ex}}

\title{Tập luận văn đội tuyển quốc gia Trung Quốc năm 2017}
\author{Tập luận văn ứng viên đội tuyển quốc gia IOI Trung Quốc}
\date{Tháng 4 năm 2017}

\begin{document}
\maketitle

\origtitle{Tập luận văn ứng viên đội tuyển quốc gia Trung Quốc Olympic Tin học năm 2017}
\sourcepath{IOI China candidate team papers / National training team 2017 collection.pdf}

\renewcommand{\abstractname}{Tóm tắt}
\begin{abstract}
PDF nguồn là tập hợp luận văn đội tuyển quốc gia Trung Quốc năm 2017. Khác với
một số tập trước đó, lớp văn bản của tập này trích xuất khá tốt bằng
\texttt{pdftotext}. Bản LaTeX này chuyển ngữ phần nhận diện tập sách, toàn bộ
mục lục, và bài đầu tiên của Mao Xiao về công thức truy hồi của dãy số.
\end{abstract}

\section{Thông tin đầu sách}

Tập sách có tiêu đề \emph{Tập luận văn ứng viên đội tuyển quốc gia Trung Quốc
Olympic Tin học năm 2017}. Trang bìa ghi thời điểm phát hành là tháng 4 năm
2017. Tệp PDF gồm 204 trang, được tạo bằng LaTeX với gói \texttt{hyperref}.

\section{Mục lục đã dịch}

\begin{center}
\small
\begin{tabular}{r p{0.52\linewidth} p{0.24\linewidth}}
\textbf{Trang} & \textbf{Bài viết} & \textbf{Tác giả} \\
\hline
1 & Một số nghiên cứu về công thức truy hồi của dãy số & Mao Xiao \\
12 & Matching trên đồ thị tổng quát dựa trên đại số tuyến tính & Yang Jiaqi \\
27 & Tính tổng đa thức & Yuan Yutao \\
32 & Bàn về bài toán tập độc lập trong thi tin học & Zhong Zhixian \\
50 & Báo cáo ra đề và mở rộng của \emph{Đồ thị con kỳ diệu} & Chen Junkun \\
75 & Tìm hiểu bài toán bao đóng bắc cầu động & Sun Yaofeng \\
90 & Báo cáo ra đề \emph{A + B Problem} & Wang Leping \\
99 & Bước đầu tìm hiểu thuật toán chia khối kích thước phi chuẩn & Xu Mingkuan \\
122 & Cây palindrome và ứng dụng & Weng Wentao \\
139 & Báo cáo ra đề \emph{Cây đen trắng} & Yan Shuyi \\
144 & Báo cáo ra đề \emph{Đa giác đều} & Yang Jingqin \\
155 & Bàn về lời giải tuyến tính cho quy hoạch động có đơn điệu quyết định & Feng Zhe \\
165 & Báo cáo ra đề \emph{Cây đoạn bị thao túng} & Shen Rui \\
179 & Bước đầu về logic máy tính và nghệ thuật: mô hình cường độ nốt khi chơi piano dựa trên logic & Zhao Shengyu \\
190 & Báo cáo ra đề \emph{Tái cấu trúc bộ gen} & Hong Huadun
\end{tabular}
\end{center}

\section{Một số nghiên cứu về công thức truy hồi của dãy số}

\begin{center}
Mao Xiao, Trường Trung học Yali
\end{center}

\renewcommand{\abstractname}{Tóm tắt}
\begin{abstract}
Bài viết này là luận văn đội tuyển quốc gia tin học Trung Quốc năm 2017 của
tác giả. Trong bài, tác giả trước hết giới thiệu đa thức truy hồi, sau đó
giới thiệu thuật toán Berlekamp--Massey, rồi trình bày triển vọng ứng dụng của
chúng trong thi tin học, gồm đếm công thức truy hồi và tìm đa thức đặc trưng
của một số ma trận thưa.
\end{abstract}

\subsection*{Dẫn nhập}

\emph{Đa thức truy hồi} là một thuật ngữ do tác giả đặt ra. Trong thi tin học,
công thức truy hồi thường được cho tường minh, còn nhiệm vụ của thí sinh
thường là: biết một số hạng đầu của dãy và một công thức truy hồi nào đó, hãy
tính một số hạng của dãy. Tuy nhiên, bài viết này nghiên cứu những công thức
truy hồi ẩn. Vì vậy tác giả đưa ra khái niệm hoàn toàn mới là đa thức truy
hồi. Đây là một khái niệm rất mạnh: nó không chỉ có các ứng dụng như đếm công
thức truy hồi, mà còn giúp ta học và hiểu thuật toán Berlekamp--Massey.

Berlekamp--Massey là một thuật toán xuất sắc nhưng bị lạnh nhạt. Trong thi tin
học, Berlekamp--Massey không mấy phổ biến; ngay cả những người biết nó cũng
thường chỉ xem nó như một đường tắt. Ngay cả trên phạm vi thế giới, tôi cũng
chưa nghe nói thuật toán này từng xuất hiện như thuật toán chuẩn của một bài
nào. Tuy vậy, nó có khá nhiều ứng dụng.

Do đó, bài viết này sẽ trình bày hai nội dung ấy và giới thiệu một số ứng dụng
của chúng.

\subsection{Đa thức truy hồi}

Trong mục này, tác giả sẽ đưa ra nhiều định nghĩa mới, đồng thời giới thiệu
nhiều kết luận hữu ích.

\subsubsection{Bậc nhỏ nhất}

\paragraph{Định nghĩa 1.1.}
Với dãy \(\{s_0,s_1,\ldots,s_{l-1}\}\), ta đặt đa thức tương ứng của dãy là
\[
  S(x)=\sum_{k=0}^{l-1}s_kx^k.
\]

\paragraph{Định nghĩa 1.2.}
Với một đa thức khác không \(A\), \term{bậc} của nó được định nghĩa là bậc của
hạng tử cao nhất, ký hiệu \(\deg(A)\). Riêng đa thức không có bậc là vô cùng
nhỏ.

\paragraph{Định nghĩa 1.3.}
Trên một trường nào đó, với dãy \(\{r_0,r_1,\ldots,r_{m-1}\}\) và dãy
\(\{a_0,a_1,\ldots,a_{n-1}\}\), nếu \(r\) không toàn bằng không và với mọi
\(0\le p\le n-m\) ta có
\[
  \sum_{k=0}^{m-1}a_{p+k}\times r_{m-k-1}=0,
  \tag{1.1}
\]
thì gọi \(r\) là một \term{công thức truy hồi}\footnote{Một số nơi gọi khái
niệm này là công thức truy tiến. Tác giả không tìm được khác biệt giữa hai
thuật ngữ, nên trong bài chọn thuật ngữ ít dùng hơn để định nghĩa riêng, còn
thuật ngữ kia giữ nghĩa thường gặp.} của dãy \(a\).

Nếu \(r_0=1\), ta gọi \(r\) là \term{công thức truy tiến}. Nếu \(r\) là công
thức truy tiến có độ dài ngắn nhất trong mọi công thức truy tiến của dãy
\(a\), ta gọi nó là \term{công thức truy tiến ngắn nhất} của dãy \(a\).

\paragraph{Định nghĩa 1.4.}
Với dãy \(a\) và đa thức khác không \(R(x)\), luôn tồn tại duy nhất công thức
truy hồi ngắn nhất \(r\) sao cho đa thức tương ứng của nó là \(R(x)\). Ta gọi
\(r\) là \term{công thức truy hồi ngắn nhất tương ứng} của \(R(x)\). Độ dài của
\(r\) trừ một được gọi là \term{bậc nhỏ nhất} của đa thức \(R(x)\), ký hiệu
\(d_R\) hoặc \(d_{R(x)}\). Đa thức \(R(x)\) được gọi là \term{đa thức truy hồi
bậc \(d_R\)} của \(a\).

\paragraph{Chứng minh.}
Cho một đa thức khác không \(R(x)\) và một dãy \(a\) độ dài \(n\). Dãy gồm các
hệ số từ hạng tử bậc \(0\) đến hạng tử bậc \(\max(\deg(R),n)\) của nó chắc
chắn là một công thức truy hồi của \(a\), và đa thức tương ứng là \(R(x)\), nên
\(r\) tồn tại.

Nếu bậc nhỏ nhất của đa thức khác không \(R(x)\) là \(d\), thì \(r\) nhất định
là dãy gồm các hệ số từ bậc \(0\) đến bậc \(d\) của nó, nên \(r\) là duy nhất.
\qedmark

\subsubsection{Các phép toán trên đa thức truy hồi}

Sau khi có đa thức truy hồi, ta có thể nhận được một số kết luận hữu ích. Từ
các kết luận này, ta có thể làm nhiều việc, và chúng cũng giúp ích cho thuật
toán Berlekamp--Massey.

\paragraph{Bổ đề 1.1.}
Đa thức \(R(x)\) có bậc nhỏ nhất \(d\) đối với dãy \(a\) khi và chỉ khi \(d\)
không nhỏ hơn bậc của \(R(x)\), và các hệ số từ bậc \(d\) đến bậc \(n-1\) của
\(R(x)\times A(x)\) đều bằng không.

\paragraph{Chứng minh.}
Ta nhận thấy công thức (1.1) có dạng tích chập; từ đó dễ suy ra bổ đề.
\qedmark

\paragraph{Định lý 1.1.}
Giả sử \(F(x),G(x)\) là hai đa thức có bậc nhỏ nhất không vượt quá \(d\) đối
với dãy \(a\). Khi đó \(F(x)+G(x)\) cũng có bậc nhỏ nhất không vượt quá \(d\)
đối với dãy \(a\).

\paragraph{Chứng minh.}
Theo Bổ đề 1.1, các hệ số từ bậc \(d\) đến bậc \(n-1\) của
\(F(x)\times A(x)\) và \(G(x)\times A(x)\) đều bằng không. Vì vậy các hệ số từ
bậc \(d\) đến bậc \(n-1\) của \((F(x)+G(x))\times A(x)\) đều bằng không; đồng
thời bậc của \(F(x)+G(x)\) hiển nhiên không vượt quá \(d\). Do đó \(F(x)+G(x)\)
có bậc nhỏ nhất không vượt quá \(d\) đối với dãy \(a\). \qedmark

\paragraph{Bổ đề 1.2.}
Giả sử \(F(x)\) có bậc nhỏ nhất \(d\) đối với dãy \(a\). Khi đó với
\(k\in\mathbb{N}\) và \(k\ge 0\), \(x^kF(x)\) có bậc nhỏ nhất không vượt quá
\(d+k\) đối với dãy \(a\).

\paragraph{Chứng minh.}
Theo Bổ đề 1.1, các hệ số từ bậc \(d\) đến bậc \(n-1\) của \(F(x)\times A(x)\)
đều bằng không. Vì vậy các hệ số từ bậc \(d+k\) đến bậc \(n+k-1\) của
\(x^kF(x)\times A(x)\) đều bằng không; đồng thời bậc của \(x^kF(x)\times A(x)\)
hiển nhiên không vượt quá \(d+k\). Do đó \(x^kF(x)\) có bậc nhỏ nhất không
vượt quá \(d+k\) đối với dãy \(a\). \qedmark

\paragraph{Định lý 1.2.}
Giả sử \(F(x)\) có bậc nhỏ nhất \(d_F\) đối với dãy \(a\). Khi đó
\(F(x)\times G(x)\) có bậc nhỏ nhất không vượt quá \(d_F+\deg(G)\) đối với dãy
\(a\).

\paragraph{Chứng minh.}
Suy ra từ Định lý 1.1 và Bổ đề 1.2. \qedmark

\subsubsection{Cơ sở}

\paragraph{Định nghĩa 1.5.}
Với dãy \(a\) và dãy đa thức
\(\{\Base_0(x),\Base_1(x),\Base_2(x),\ldots\}\), nếu với mọi \(k\),
\(\Base_k(x)\) là đa thức có bậc nhỏ nhất nhỏ nhất trong mọi đa thức có hạng
tử thấp nhất là \(x^k\), thì gọi dãy đa thức \(\Base\) là \term{một cơ sở} của
dãy \(a\).

Ta nhận thấy một cơ sở của một dãy thật ra chính là cơ sở của không gian tuyến
tính do nó sinh ra. Ta cũng có thể nhận thấy tọa độ của một đa thức theo một
cơ sở liên hệ chặt chẽ với bậc nhỏ nhất của nó.

\paragraph{Định lý 1.3.}
Nếu \(\Base\) là cơ sở của dãy \(a\), thì mọi đa thức có bậc nhỏ nhất \(d\) đều
có thể được biểu diễn tuyến tính một cách duy nhất bởi một số \(\Base_k(x)\) có
bậc nhỏ nhất không vượt quá \(d\); hơn nữa trong biểu diễn đó nhất định tồn
tại ít nhất một \(\Base_k(x)\) có bậc nhỏ nhất đúng bằng \(d\).

\paragraph{Chứng minh.}
Giả sử \(R(x)\) là một đa thức có bậc nhỏ nhất \(d\). Ta xét cách phân rã nó.

Nếu \(R(x)\) bằng không thì không cần phân rã. Ngược lại, giả sử hạng tử thấp
nhất của nó là \(x^k\), hệ số là \(C_r\), và hệ số tương ứng của \(\Base_k(x)\)
là \(C_{\Base}\). Khi đó hạng tử thấp nhất của
\(R(x)-C_{\Base}^{-1}C_r\Base_k(x)\) ít nhất là \(x^{k+1}\). Vì \(\Base_k(x)\)
là một đa thức có bậc nhỏ nhất nhỏ nhất trong các đa thức có hạng tử thấp nhất
\(x^k\), bậc nhỏ nhất của nó hiển nhiên không vượt quá \(d\). Theo Định lý 1.1,
\(R(x)-C_{\Base}^{-1}C_r\Base_k(x)\) cũng có bậc nhỏ nhất không vượt quá \(d\).
Ta phân rã biểu thức này theo cách đệ quy. Cuối cùng ta thu được một biểu diễn
tuyến tính thỏa điều kiện.

Vì các \(\Base_k(x)\) là một cơ sở của không gian tuyến tính mà chúng sinh ra,
cách biểu diễn hiển nhiên là duy nhất.

Nếu trong biểu diễn, mọi \(\Base_k(x)\) đều có bậc nhỏ nhất nhỏ hơn \(d\), thì
hoặc đa thức này là đa thức không, hoặc theo Bổ đề 1.1 bậc nhỏ nhất của nó sẽ
nhỏ hơn \(d\), mâu thuẫn. \qedmark

\subsection{Thuật toán Berlekamp--Massey}

Mục này giới thiệu thuật toán Berlekamp--Massey\footnote{Wikipedia tiếng Trung
phiên âm thuật toán này là Berlekamp--Massey.}. Tuy các tài liệu như Wikipedia
\([1]\) đều đã giới thiệu quy trình của thuật toán Berlekamp--Massey, tác giả
không tìm được bất kỳ bản quy trình thuật toán tiếng Trung nào. Đồng thời, tác
giả cũng muốn dùng ngôn ngữ của bài viết này để trình bày lại quy trình của
thuật toán.

\subsubsection{Giới thiệu thuật toán}

Thuật toán Berlekamp--Massey có thể tìm một cơ sở \(\Base\) cho dãy \(a\) độ
dài \(n\) trên một trường bất kỳ bằng \(\Oh(n^2)\) phép toán và không cần thêm
không gian đáng kể.

\subsubsection{Quy trình thuật toán}

Theo định nghĩa, mọi đa thức có hạng tử thấp nhất \(x^k\), với \(k\ge n\), đều
có bậc nhỏ nhất đúng bằng bậc của nó. Vì vậy với \(\Base_k(x)\), ta chỉ cần
chọn bất kỳ \(Cx^k\) nào với \(C\ne 0\). Không mất tính tổng quát, giả sử mọi
đa thức tìm được đều có hệ số của hạng tử thấp nhất bằng \(1\), khi đó ở đây
ta đặt \(C=1\).

Giả sử ta đã tìm được một đa thức có bậc nhỏ nhất nhỏ nhất trong các đa thức
từ \(\Base_{k+1}(x)\) đến \(\Base_n(x)\). Xét cách tìm \(\Base_k(x)\).

Ta xét \(x\Base_k(x)\). Theo Định lý 1.3, ta chỉ cần làm cho bậc nhỏ nhất lớn
nhất trong biểu diễn tuyến tính của nó càng nhỏ càng tốt.

Trước hết, vì hạng tử thấp nhất của \(x\Base_k(x)\) là \(x^{k+1}\), trong biểu
diễn nhất định có hạng \(\Base_{k+1}(x)\). Xét các hạng còn lại: rõ ràng bậc
nhỏ nhất lớn nhất trong số chúng phải càng nhỏ càng tốt.

Đặt \(\Res(x)=x\Base_k(x)-\Base_{k+1}(x)\). Ta tính
\(\Pro(x)=\Res(x)\times A(x)\). Theo Bổ đề 1.1, dễ biết các hệ số từ bậc
\(k+1\) đến bậc \(n-1\) của \(\Pro(x)\) đều bằng không.

Xét hạng tử bậc \(n\). Nếu nó bằng không, ta đã biết \(x\Base_k(x)\); ngược lại
giả sử nó là \(u\). Xét \(\Base_l(x)\) với \(l>k+1\). Nếu hệ số bậc \(n\) của
\(\Base_l(x)\times A(x)\) khác không, giả sử hệ số bậc \(n\) là \(v\). Dễ biết
các hệ số từ bậc \(k+1\) đến bậc \(n\) của
\(\Res(x)-uv^{-1}\Base_l(x)\) đều bằng không. Vì bậc nhỏ nhất lớn nhất phải
càng nhỏ càng tốt, nếu tồn tại đa thức thỏa điều kiện thì \(\Base_l(x)\) nhất
định là đa thức có bậc nhỏ nhất nhỏ nhất trong mọi đa thức thỏa điều kiện. Lấy
\(\Base_l(x)\) như vậy, ta có thể tìm được \(x\Base_k(x)\). Nếu không tồn tại
\(\Base_l(x)\) như vậy, thì chỉ có thể là bậc nhỏ nhất của \(x\Base_k(x)\) đã
vượt quá \(n\). Khi đó ta chỉ cần tùy ý chọn một đa thức có hạng tử thấp nhất
là hạng bậc \(x^{k+1}\) và hệ số bằng \(1\).

Khi đã biết \(x\Base_k(x)\), hiển nhiên ta cũng biết \(\Base_k(x)\).

Dùng phương pháp này, mỗi \(P_k(x)\) đều có thể được tìm trong không quá
\(\Oh(n)\) phép toán, vì vậy tổng số phép toán là \(\Oh(n)\), còn số \(\Base\)
cần tính chỉ là \(\Oh(n)\). Do đó tổng số phép toán là \(\Oh(n^2)\).

Rõ ràng ngoài việc lưu \(A(x)\) và \(\Base\), thuật toán này không cần không
gian phụ đáng kể.

\subsubsection{Mã giả}

Theo mô tả trên, thuật toán có thể được cài đặt bằng phương pháp sau.

\begin{center}
\small
\begin{tabular}{r@{\quad}p{0.78\linewidth}}
\multicolumn{2}{l}{\textbf{Thuật toán 1} Berlekamp--Massey Algorithm I}\\
\multicolumn{2}{l}{\textbf{Require:} \(a, n=|a|\)}\\
\multicolumn{2}{l}{\textbf{Ensure:} \(\Base\)}\\
1: & \(\Base_k=x^k\ (k\ge n)\)\\
2: & \(l\gets n+1\)\\
3: & \(v\gets 1\)\\
4: & \textbf{for} \(k\gets n-1\) \textbf{to} \(1\) \textbf{do}\\
5: & \quad \(u\gets [x^n]\Base_{k+1}(x)\times A(x)\)\\
6: & \quad \(\Base_k(x)\gets(\Base_{k+1}(x)-uv^{-1}\Base_l(x))\div x\)\\
7: & \quad \textbf{if} \((u\ne 0)\wedge(d_{\Base_{k+1}}<d_{\Base_l})\)
      \textbf{then}\\
8: & \qquad \(l\gets k+1\)\\
9: & \qquad \(v\gets u\)\\
10: & \quad \textbf{end if}\\
11: & \textbf{end for}=0
\end{tabular}
\end{center}

Theo mô tả thuật toán, giá trị ban đầu của \(l\) và \(v\) không quan trọng;
chỉ cần giá trị ban đầu thỏa \(l\ge n+1\) là có thể bảo đảm đáp án đúng. Tuy
nhiên, \(l=n+1, v=1\) là giá trị khởi tạo tham khảo trên Wikipedia tiếng Anh,
nên tác giả xem nó như một phương án đã thành thông lệ.

Tuy nhiên, vì phép toán đa thức khá phiền phức, khi cài đặt thật sự ta thường
không dùng đa thức một cách tường minh.

Ký hiệu \(r_k\) là công thức truy hồi ngắn nhất tương ứng với
\(\Base_{n-k+1}\div x^{n-k+1}\). Khi đó ta dễ thấy \(r_k\) là công thức truy
tiến ngắn nhất của tiền tố độ dài \(k\) của dãy \(a\). Vì vậy quá trình trên
dễ dàng được đổi thành quá trình tính \(r\).

Ta xem dãy như một mảng có độ dài thay đổi, chỉ số bắt đầu từ \(0\). Thuật toán
mới như sau.

\begin{center}
\small
\begin{tabular}{r@{\quad}p{0.78\linewidth}}
\multicolumn{2}{l}{\textbf{Thuật toán 2} Berlekamp--Massey Algorithm II}\\
\multicolumn{2}{l}{\textbf{Require:} \(a, n=|a|\)}\\
\multicolumn{2}{l}{\textbf{Ensure:} \(r_k\ (0\le k\le n)\)}\\
1: & \(r_0=\{1\}\)\\
2: & \(\last\gets\{1\}\)\\
3: & \(v\gets 1\)\\
4: & \(s\gets 1\)\\
5: & \textbf{for} \(k\gets 1\) \textbf{to} \(n\) \textbf{do}\\
6: & \quad \(r_k\gets r_{k-1}\)\\
7: & \quad \(u\gets\sum_{i=0}^{|r_k|-1}a_{k-i-1}\times r_i\)\\
8: & \quad \textbf{if} \((u\ne 0)\) \textbf{then}\\
9: & \qquad \textbf{if} \(|r_k|<|\last|+s\) \textbf{then}\\
10: & \qquad\quad \(r_k.\append(\{0\}\times(|\last|+s-|r_k|))\)\\
11: & \qquad\quad \textbf{for} \(i\gets 0\) \textbf{to} \(|\last|-1\)
       \textbf{do}\\
12: & \qquad\qquad \(r_k[i+s]\gets r_k[i+s]-uv^{-1}\last[i]\)\\
13: & \qquad\quad \textbf{end for}\\
14: & \qquad\quad \(\last\gets r_{k+1}\)\\
15: & \qquad\quad \(v\gets u\)\\
16: & \qquad\quad \(s\gets 0\)\\
17: & \qquad \textbf{else}\\
18: & \qquad\quad \textbf{for} \(i\gets 0\) \textbf{to} \(|\last|-1\)
       \textbf{do}\\
19: & \qquad\qquad \(r_k[i+s]\gets r_k[i+s]-uv^{-1}\last[i]\)\\
20: & \qquad\quad \textbf{end for}\\
21: & \qquad \textbf{end if}\\
22: & \quad \textbf{end if}\\
23: & \quad \(s\gets s+1\)\\
24: & \textbf{end for}=0
\end{tabular}
\end{center}

Ngoài ra, nếu ta chỉ cần tìm công thức truy tiến ngắn nhất của cả dãy, thì ở
mọi thời điểm ta thật ra chỉ cần lưu một \(r_k\) và một \(\last\). Nếu không
tính kích thước từng phần tử, độ phức tạp không gian có thể giảm xuống
\(\Oh(n)\).

\subsection{Đếm công thức truy hồi}

\subsubsection{Đa thức truy hồi nhỏ nhất và đa thức truy hồi thứ nhì không tầm thường}

\paragraph{Định nghĩa 3.1.}
Trong mọi \(\Base_k(x)\) đã tìm được, một đa thức có bậc nhỏ nhất nhỏ nhất
được gọi là \term{đa thức truy hồi nhỏ nhất}, ký hiệu \(M(x)\). Nếu có nhiều
đa thức như vậy thì chọn tùy ý; bậc nhỏ nhất của nó được ký hiệu \(d_M\).

Trong mọi \(\Base_k(x)\) đã tìm được, nếu tồn tại một đa thức không thể biểu
diễn dưới dạng \(x^kM(x)\), với \(k\in\mathbb{N}, k\ge 0\), thì đa thức có bậc
nhỏ nhất nhỏ nhất trong số đó được gọi là \term{đa thức truy hồi thứ nhì không
tầm thường}, ký hiệu \(S(x)\). Nếu có nhiều đa thức như vậy thì chọn tùy ý; nếu
không tồn tại thì đặt \(S(x)=\Base_{n+1}(x)\). Bậc nhỏ nhất của nó được ký hiệu
\(d_S\).

\subsubsection{Biểu diễn tuyến tính}

\paragraph{Định nghĩa 3.2.}
Với các đa thức \(A(x),B(x),C(x)\), nếu tồn tại duy nhất đa thức \(F_B(x)\) có
bậc không vượt quá \(d_A-d_B\) và đa thức \(F_C(x)\) có bậc không vượt quá
\(d_A-d_C\) sao cho
\[
  A(x)=F_B(x)B(x)+F_C(x)C(x),
\]
thì gọi \(A(x)\) là \term{có thể biểu diễn tuyến tính} bởi \(B(x),C(x)\).

\paragraph{Bổ đề 3.1.}
Với các đa thức \(A(x),B(x),C(x),D(x),E(x)\), nếu \(A(x)\) có thể biểu diễn
tuyến tính bởi \(B(x),C(x)\), và \(B(x),C(x)\) đều có thể biểu diễn tuyến tính
bởi \(D(x),E(x)\), thì \(A(x)\) cũng có thể biểu diễn tuyến tính bởi
\(D(x),E(x)\).

\paragraph{Chứng minh.}
Đặt
\[
\begin{aligned}
  A(x)&=F_B(x)B(x)+F_C(x)C(x),\\
  B(x)&=F_{BD}(x)D(x)+F_{BE}(x)E(x),\\
  C(x)&=F_{CD}(x)D(x)+F_{CE}(x)E(x).
\end{aligned}
\]
Dễ biết
\[
\begin{aligned}
  A(x)={}&(F_B(x)F_{BD}(x)+F_C(x)F_{CD}(x))D(x)\\
        &+(F_D(x)F_{BE}(x)+F_C(x)F_{CE}(x))E(x).
\end{aligned}
\]
Vì
\[
  \deg(F_B(x))\le d_A-d_B,\qquad
  \deg(F_{BD}(x))\le d_B-d_D,
\]
nên
\[
  \deg(F_B(x)F_{BD}(x))
  =\deg(F_B(x))+\deg(F_{BD}(x))
  \le d_A-d_D.
\]
Tương tự,
\[
  \deg(F_C(x)F_{CD}(x))\le d_A-d_D.
\]
Do đó
\[
  \deg(F_B(x)F_{BD}(x)+F_C(x)F_{CD}(x))\le d_A-d_D.
\]
Tương tự,
\[
  \deg(F_B(x)F_{BE}(x)+F_C(x)F_{CE}(x))\le d_A-d_E.
\]
Quá trình này bảo đảm tính duy nhất, nên bổ đề được chứng minh. \qedmark

\paragraph{Định lý 3.1.}
Sau khi tính xong \(\Base_k(x)\), nó là đa thức nhỏ nhất hoặc là đa thức thứ
nhì không tầm thường.

Tạm thời ta chưa chứng minh định lý này; trước hết giới thiệu một bổ đề mới,
rồi chứng minh cả hai cùng lúc.

\paragraph{Bổ đề 3.2.}
Mọi \(\Base_k(x)\) đều có thể được biểu diễn tuyến tính bởi \(S(x)\) và
\(M(x)\).

\paragraph{Chứng minh.}
Xét quy trình của thuật toán Berlekamp--Massey.

Khi \(k=n\), hiển nhiên Định lý 3.1 và Bổ đề 3.2 đều đúng.

Giả sử sau khi tính xong \(\Base_{k+1}(x)\), cả hai đều đúng.

Xét quá trình tính \(x\Base_k(x)\). Cuối cùng nó nhất định là
\(\Base_{k+1}(x)\) cộng với một hằng số nhân với một \(\Base_l(x)\) nào đó.
Theo mô tả thuật toán, ta dễ nhận thấy \(\Base_{k+1}(x)\) và hạng còn lại đều
là một trong \(M(x)\) và \(S(x)\).

Vì vậy bậc nhỏ nhất của \(x\Base_k(x)\) bằng \(\max(d_M,d_S)\), còn bậc nhỏ
nhất của \(\Base_k(x)\) phải trừ thêm một. Do đó nó nhất định là đa thức nhỏ
nhất hoặc đa thức thứ nhì tầm thường. Như vậy ta đã chứng minh Định lý 3.1.

Hiển nhiên với chính nó, Bổ đề 3.2 đúng. Tuy nhiên vì \(M(x),S(x)\) đã thay
đổi, ta còn phải xét \(\Base_l(x)\) với \(l>k\).

Xét sự thay đổi của \(M(x)\) và \(S(x)\). \(M(x)\) hoặc không đổi, hoặc trở
thành \(S(x)\) mới.

Xét \(S(x)\). Theo quy trình thuật toán, ta có
\[
  x\Base_k(x)=C_MM(x)+C_SS(x).
\]
Do đó
\[
  S(x)=x\Base_k(x)-C_MM(x).
\]
Vì \(d_S>d_{\Base_k}\) và \(d_S\ge d_M\), \(S(x)\) trước khi thay đổi có thể
được biểu diễn tuyến tính bởi \(M(x)\) và \(S(x)\) sau khi thay đổi.

Như vậy, \(M(x)\) và \(S(x)\) trước khi thay đổi đều có thể được biểu diễn
tuyến tính bởi \(M(x)\) và \(S(x)\) sau khi thay đổi. Theo Bổ đề 3.1,
\(\Base_l(x)\) với \(l>k\) vẫn có thể được biểu diễn tuyến tính bởi \(S(x)\) và
\(M(x)\).

Đồng thời, với trường hợp \(k>n\), vì \(\Base_k(x)=x^{k-n}\Base_n(x)\), bổ đề
hiển nhiên đúng.

Tương tự, quá trình này giữ tính duy nhất, nên bổ đề được chứng minh.
\qedmark

\paragraph{Định lý 3.2.}
Mọi đa thức đều có thể được biểu diễn tuyến tính bởi \(S(x)\) và \(M(x)\).

\paragraph{Chứng minh.}
Dễ suy ra từ Định lý 1.3 và Bổ đề 3.2. \qedmark

\subsubsection{Định lý đếm công thức truy hồi}

\paragraph{Định lý 3.3 (Định lý đếm đa thức truy hồi).}
Giả sử kích thước của trường là \(q\). Khi đó số đa thức truy hồi bậc \(k\) của
một dãy là
\[
  \num_k=
  q^{\max(k-d_M-1,0)+\max(k-d_S-1,0)}
  \left(q^{[k\ge d_M]+[k\ge d_S]}-1\right).
\]
Trong công thức này, \([x]\) bằng \(1\) khi \(x\) đúng, và bằng \(0\) nếu
không.

\paragraph{Chứng minh.}
Dễ suy ra từ Định lý 3.2. \qedmark

\paragraph{Định lý 3.4 (Định lý đếm công thức truy hồi).}
Số công thức truy hồi của một dãy là
\[
  \sum_{k=1}^{n}(n-k+1)\num_k.
\]

\paragraph{Chứng minh.}
Một đa thức truy hồi bậc \(x\) tương ứng đúng một công thức truy hồi có độ dài
\(x,x+1,\ldots,n\), tổng cộng \(n-x+1\) công thức, nên có công thức trên.
\qedmark

\subsection{Đa thức đặc trưng của ma trận thưa đặc biệt}

Thuật toán Berlekamp--Massey không chỉ có tác dụng rất mạnh trong lĩnh vực
công thức truy hồi; nó cũng có thể giúp ích cho các bài toán liên quan đến ma
trận.

Vì vậy, chương này sẽ giới thiệu cách dùng thuật toán đó để tính đa thức đặc
trưng của ma trận thưa đặc biệt.

\subsubsection{Điều kiện quan trọng}

Khi mỗi trị riêng chỉ tương ứng với một khối Jordan\footnote{Nói cách khác,
chiều cao của mỗi vector đều bằng bội số của nó.}, đa thức đặc trưng chính là
đa thức tối tiểu. Đây là điều kiện quan trọng để thuật toán này hoạt động.

\subsubsection{Quá trình thực hiện}

Theo định lý Cayley--Hamilton, đa thức đặc trưng của một ma trận \(n\times n\)
\(A\) là đa thức triệt tiêu của nó; tức là nếu đa thức đặc trưng là \(p\) thì
\(p(A)=0\).

Nhớ lại định nghĩa công thức truy hồi, ta thấy nếu chọn ngẫu nhiên một vector
\(v\), trước hết hiển nhiên \(p(Av)=0\). Nếu kích thước của trường số đủ lớn,
thì với xác suất rất cao, dãy vector \(\{v,Av,A^2v,\ldots\}\) tồn tại một công
thức truy tiến ngắn nhất duy nhất \(r\) có độ dài \(n\). Khi đó \(R(x)\) chính
là đa thức đặc trưng.

Nếu \(A\) là ma trận thưa và số phần tử là \(e\), thì từ \(A^kv\) có thể tính
\(A^{k+1}v\) trong \(\Oh(e)\) phép toán.

Khi xét dùng thuật toán Berlekamp--Massey, ta có hai khó khăn.

Khó khăn thứ nhất là xử lý nhanh dãy vector. Ta có thể thiết kế một hàm băm,
trong đó hàm băm là một biến đổi tuyến tính, biến mọi vector thành một số; như
vậy ta có thể xử lý dãy vector.

Khó khăn thứ hai là xử lý một dãy vô hạn. Tuy dãy này dài vô hạn, ta có thể
nhận thấy trong trường hợp ngẫu nhiên, độ dài công thức truy tiến ngắn nhất của
tiền tố độ dài \(2n\) xấp xỉ \(n\). Vì vậy ta có thể lấy một tiền tố đủ dài để
tính. Một cách tốt hơn là xét Thuật toán 2: thật ra ta có thể liên tục tăng
tiền tố cho đến khi độ dài công thức truy tiến ngắn nhất đạt \(n\).

Như vậy, ta thu được một thuật toán Monte Carlo có thể tính đa thức đặc trưng
của ma trận thưa đặc biệt bằng \(\Oh(n(n+e))\) phép toán và không gian tuyến
tính theo kích thước ma trận.

\subsubsection{Định thức của ma trận thưa}

Xét cách tính định thức \(\det(A)\). Hiển nhiên nếu ma trận này thỏa điều kiện
quan trọng ở Mục 4.1, ta có thể tìm đa thức đặc trưng, rồi lấy hạng tử tự do
nhân với \((-1)^n\).

Nếu ma trận này không thỏa điều kiện đó, ta có thể lấy ngẫu nhiên một ma trận
đường chéo \(B\); số phép toán là \(\Oh(ne)\). Nếu trường số đủ lớn, thì \(AB\)
có xác suất rất cao thỏa điều kiện này. Vì
\(\det(AB)=\det(A)\det(B)\), ta chỉ cần tính \(\det(AB)\), rồi chia cho
\(\det(B)\). Vì \(B\) là ma trận đường chéo, \(\det(B)\) chính là tích các phần
tử trên đường chéo của \(B\).

Như vậy, ta thu được một thuật toán Monte Carlo có thể tính định thức của ma
trận thưa bằng \(\Oh(n(n+e))\) phép toán và không gian tuyến tính theo kích
thước ma trận.

\subsubsection{Đếm cây khung của đồ thị thưa}

Hiển nhiên, có thể dùng định lý ma trận-cây Kirchhoff để tính số cây khung.
Nếu đồ thị đủ thưa, thì định thức của ma trận Laplace hiển nhiên có thể được
tính bằng phương pháp ở Mục 4.3.

Độ phức tạp thời gian của thuật toán này là \(\Oh(n(n+m))\), độ phức tạp không
gian là \(\Oh(n+m)\), rất tốt.

\subsection{Tổng kết}

Toàn bài đều xoay quanh công thức truy hồi. Phần thứ nhất trình bày đa thức
truy hồi, phần thứ hai nghiên cứu thuật toán Berlekamp--Massey, còn phần thứ
ba và thứ tư chủ yếu giới thiệu một số ứng dụng của chúng.

Tôi hy vọng bài viết này có thể đóng vai trò gợi mở. Đây vẫn chỉ là một nghiên
cứu chưa chín muồi về công thức truy hồi, và các ứng dụng được giới thiệu cũng
khá cơ bản. Khi khoa học kỹ thuật không ngừng tiến bộ, nhất định còn nhiều ứng
dụng thú vị hơn đang chờ chúng ta khám phá. Vì vậy, mong rằng bài viết này có
thể gợi ý cho độc giả và dẫn ra một số ứng dụng hấp dẫn hơn.

\subsection*{Lời cảm ơn}

\begin{enumerate}
  \item Cảm ơn bạn Du Yuhao. Lần đầu tiên tôi nghe nói về thuật toán này là từ
  bạn ấy, và sự xuất hiện của một phần nội dung trong bài viết này không thể
  tách rời sự giúp đỡ của bạn ấy.
  \item Cảm ơn các bạn Du Yuhao, Yang Jiaqi và Weng Wentao đã sửa lỗi.
  \item Cảm ơn China Computer Federation đã cung cấp nền tảng học tập và trao
  đổi.
\end{enumerate}

\subsection*{Tài liệu tham khảo}

\begin{enumerate}[label={[\arabic*]}]
  \item Wikipedia.

  \url{https://en.wikipedia.org/wiki/Berlekamp%E2%80%93Massey_algorithm}
\end{enumerate}

\section*{Ghi chú về trích xuất}

Nội dung bài đầu được trích bằng \texttt{pdftotext} từ trang vật lý 3 đến 13
của PDF, tương ứng trang bài viết 1 đến 11; trang vật lý 14 là bài kế tiếp.
Không cần OCR. Hai dòng kết thúc vòng lặp trong mã giả của nguồn in là
\(\text{\textbf{end for}}=0\); bản dịch giữ nguyên ký hiệu này như một dị biệt
của nguồn.

\end{document}
